% 第一章 绪论
% 说明：本文引用采用 \cite{} 形式，文献 key 可在 bib 文件中统一维护。

\chapter{绪论}
\label{chap:introduction}

\section{研究背景与意义}
\label{sec:research-background}

随着移动互联网、智能终端和网络带宽的持续发展，在线视频已经成为互联网内容传播的重要形式。用户不再只是被动观看视频，也会通过投稿、评论、弹幕、点赞、收藏、投币和关注等方式参与内容生产与社区互动。因此，一个完整的视频分享平台不仅要支持视频播放，还需要具备内容生产、资源处理、用户互动、内容检索和后台治理等能力。

在线视频平台的业务链路较长，涉及用户管理、视频投稿、文件上传、视频转码、内容审核、搜索播放和互动统计等多个环节。若所有功能集中在单体系统中，随着功能增加，容易出现模块耦合度高、维护困难和扩展不便等问题。微服务架构强调围绕业务能力拆分服务，各服务可以独立开发、部署和扩展，适合用于功能模块较多、业务边界较清晰的 Web 平台系统\cite{lewis2014microservices,newman2021building,richardson2018microservices}。

与此同时，用户生成内容数量不断增加，平台内容安全问题也更加突出。仅依赖人工审核会带来效率低、工作量大和响应不及时等问题。语音识别、声音事件检测和多模态模型的发展，为视频内容审核提供了新的技术手段。将 AI 初审结果提供给管理员参考，可以在保证人工终审的基础上提高审核效率，形成“AI 初审 + 人工终审”的协同审核流程\cite{radford2023whisper,gemmeke2017audioset,bai2023qwen}。

基于上述背景，本文设计并实现一个基于微服务架构的在线视频分享平台。系统围绕普通用户、内容创作者和管理员三类主要对象展开，完成视频浏览、搜索、投稿、转码、互动、后台审核和 AI 辅助审核等功能。课题的实践意义在于：一方面，通过微服务拆分提高系统结构清晰度和可维护性；另一方面，通过视频处理、异步任务和智能审核等模块，构建较完整的视频平台业务闭环。

\section{国内外研究现状}
\label{sec:research-status}

国外在线视频平台起步较早，YouTube、Netflix 等平台在视频分发、流媒体播放、内容推荐和系统架构方面积累了较多实践经验。随着云计算和微服务技术的发展，大型内容平台普遍通过服务拆分、负载均衡、弹性扩展和自动化部署支撑复杂业务。与此同时，HLS 等流媒体协议和 FFmpeg 等开源音视频工具也被广泛应用于视频转码、切片和格式转换场景\cite{ffmpegdocs,rfc8216hls}。

国内在线视频和短视频平台发展迅速，哔哩哔哩、抖音、快手等平台在内容创作、弹幕互动、社区运营和内容审核方面形成了较成熟的业务模式。国内企业在后端架构上也大量采用 Spring Cloud、Nacos、Redis、Elasticsearch、RabbitMQ 等技术，以应对业务模块复杂、数据读写频繁和服务扩展困难等问题\cite{springclouddocs,nacosdocs,redisdocs,elasticsearchdocs,rabbitmqdocs}。

在内容审核方面，国内外平台均逐渐从单纯人工审核转向“机器审核 + 人工复核”的模式。传统审核方式主要依靠关键词、规则和简单分类模型，面对复杂语义和音视频混合内容时存在局限。近年来，多模态人工智能技术能够综合分析语音文本、声音事件和画面内容，为视频审核提供更丰富的判断依据\cite{radford2023whisper,tensorflow2024yamnet,bai2023qwen}。

总体来看，在线视频平台相关技术已经较为成熟，但对于毕业设计而言，难点在于如何将微服务架构、视频处理、用户互动、搜索功能和 AI 审核服务整合为一个可运行的系统。本文在现有技术基础上，构建一个具备核心业务闭环的视频分享平台原型，以体现相关技术在复杂 Web 系统中的综合应用价值。

\section{研究内容与目标}
\label{sec:research-content-goal}

本文围绕“基于微服务架构的在线视频分享平台设计与实现”展开，重点研究如何将视频平台中的用户管理、视频投稿、资源处理、互动行为、后台审核和 AI 内容审核等功能进行合理拆分与协同实现。系统采用前后端分离架构，后端以 Spring Cloud Alibaba 微服务体系为基础，AI 审核部分采用独立 Python 服务实现。

本文主要研究内容如下：

\begin{enumerate}
  \item 完成在线视频分享平台的需求分析与总体设计。根据平台业务特点，明确普通用户、内容创作者、管理员和 AI 审核服务的功能边界，并设计系统总体架构和服务模块划分。
  \item 实现视频投稿、上传、转码和播放链路。系统支持创作者上传视频和封面，保存投稿信息，并通过资源服务完成视频合并、转码和 HLS 播放资源生成\cite{ffmpegdocs,rfc8216hls}。
  \item 实现视频浏览、搜索和互动功能。用户可以浏览公开视频、按关键词检索视频、进入详情页播放视频，并进行评论、弹幕、点赞、收藏、投币和关注等互动操作。
  \item 实现后台管理与人工审核功能。管理员可以维护视频分类、查看待审核稿件、查看 AI 审核结果，并对视频进行通过或驳回处理。
  \item 实现独立 AI 多模态审核服务。AI 服务通过消息队列接收审核任务，对视频语音、声音事件和关键帧内容进行分析，并将风险等级、风险标签和审核建议返回后端\cite{rabbitmqdocs,radford2023whisper,bai2023qwen,fastapidocs}。
\end{enumerate}

本文的总体目标是完成一个功能相对完整、结构清晰、流程闭环的在线视频分享平台原型。系统应能够体现微服务架构在复杂业务拆分中的应用，能够完成视频从上传、处理、审核到发布播放的主要流程，并能够通过 AI 审核服务提高内容审核的辅助决策能力。

\section{论文结构安排}
\label{sec:thesis-structure}

本文共分为七章，各章内容安排如下：

\begin{table}[htbp]
  \centering
  \caption{论文章节结构安排}
  \label{tab:thesis-structure}
  \begin{tabular}{p{0.18\textwidth}p{0.72\textwidth}}
    \hline
    \textbf{章节} & \textbf{主要内容} \\
    \hline
    第一章 & 介绍课题研究背景、研究现状、研究内容与论文结构。 \\
    第二章 & 介绍系统开发涉及的关键技术，包括微服务、数据存储、视频处理、消息队列和 AI 审核技术。 \\
    第三章 & 从角色、功能、非功能需求和可行性等方面分析系统需求。 \\
    第四章 & 说明系统总体架构、服务划分、数据库、缓存搜索、文件存储和核心流程设计。 \\
    第五章 & 介绍用户端、创作中心、资源处理、后台管理和 AI 审核等模块的详细实现。 \\
    第六章 & 对系统核心功能、视频处理流程、AI 审核流程和异常场景进行测试与结果分析。 \\
    第七章 & 总结本文完成的工作，分析项目不足，并提出后续优化方向。 \\
    \hline
  \end{tabular}
\end{table}
